\paragraph{The quotient by vertex permutations.}
Let \(S_3\) permute the vertices \(0,1,3\), fixing vertex 2. This action
preserves the pure curve \(C\), since it permutes rows and columns of
the permanent simultaneously and preserves the vertex-2 sum.
It lifts to the normalization \(\widetilde C\).
The smooth projective model of the invariant plane curve in \((u,v)\)
is \(\widetilde C/S_3\) and has genus seven.

Indeed, work on the dense open with distinct incident edges, normalized
\(w=1\), and \(qu(v-3)\ne0\). The excluded edge and interpolation
sections above are finite. The projective direction there is recovered
from \(p,q,u,v\) by ordering the three roots of \(X^3+pX-q\) and
setting \(y_i=x_i^2+vx_i+u\). After eliminating \(q\), the first
remaining equation is quartic in \(p\) with leading coefficient 4.
Thus the map to \((u,v)\) has finite fibers and dense image on the
irreducible plane curve. Its first subresultant recovers \(p\)
generically, then the quadratic equation recovers \(q\). The splitting
field of the incident cubic therefore gives \(k(C)\) over \(k(u,v)\),
of degree at most six. All six permutations act faithfully on the
ordered roots and fix \(u,v\), proving \(k(C)^{S_3}=k(u,v)\).

For the ramification, work over \(\bar k\). Use incident and opposite
triples \(x=(x_0,x_1,x_3)\) and \(y=(y_0,y_1,y_3)\). A projective
fixed point of a transposition must be in one of its two eigenspaces.
For \((01)\), the incident triples scale respectively to \((1,1,-2)\)
and \((1,-1,0)\), since the all-zero incident triple is excluded.
The corresponding opposite triples are \((h,h,j)\) and \((h,-h,0)\).
In both cases the three pure equations generate the unit ideal.
For a three-cycle, the invariant eigenvalue would force \(x=0\).
The other two eigenspaces have \(x=(1,\omega,\omega^2)\),
\(y=vx\), for the two primitive cube roots \(\omega\).
Substitution gives \(\omega^2+\omega+1=0\) and \(v^2(v+2)=0\).
The two points with \(v=0\) are the singularities computed above.
The three-cycle fixes each of their two branches, since its order is
three. The two points with \(v=3\) are smooth. This gives six fixed
points on \(\widetilde C\), each with stabilizer of order three, and
three geometric branch values on the quotient. All ramification is
tame in characteristic five. Riemann--Hurwitz \cite{stacksTameHurwitz} gives
\[
       2\cdot43-2=6(2g-2)+6(3-1),\qquad g=7.
\]
All inertia lies in \(A_3\), so the intermediate double cover
\(\widetilde C/A_3\to\widetilde C/S_3\) is everywhere unramified
and has source genus 13. Its function field adjoins the square root of
\(-4p^3-27q^2\), since the Vandermonde product in the ordered roots
changes sign under each odd permutation and squares to the discriminant.

\paragraph{All three relabelings of the pencil.}
Write \(Y=\sum_i y_i\), \(z_i=x_i-y_i\). The three distinct pencils
are \([z_i:m_i]\) with \(m_i=Y+z_i\). The original pencil has
\(i=3\). Every pencil excludes 136 geometric directions on \(C\)
when its section at infinity is included. At \(Y=0\) all pencil values
are 1, giving the 16 geometric points of \(d=a+m\).
Elsewhere normalize \(Y=1\). Then \(\sum_i z_i=-1\), and \([z_i:1+z_i]\)
is rational over \(\mathbb F_5\) exactly when \(z_i\) is rational.
Two rational entries force all three to be rational. For the 25
ordered triples in \(\mathbb F_5^3\) of sum \(-1\), substitution of
\(y_i=x_i-z_i\) in the pure equations gives geometric points only
for permutations of \((1,1,2)\) and \((1,4,4)\), respectively two and
six per ordering. The triple intersection thus has \(16+6+18=40\)
distinct points, and every pairwise intersection equals it. The three
pencils together exclude \(3\cdot136-2\cdot40=328\) geometric
directions. These counts concern points on the original pure curve.
For the finite sections alone, the pattern \((1,4,4)\) contributes
two values at infinity. The finite triple overlap has 22 points and
the finite union has \(3\cdot112-2\cdot22=292\).

This exclusion can be evaluated on the quotient. In the \(w=1\) chart
put \(Y=3u-2p\) and
\[
 H(T)=\operatorname{Res}_X\!\left(X^3+pX-q,\,
             T+X^2-(1-v)X+u\right)=\prod_i(T-z_i).
\]
Every surviving direction satisfies \(Y\ne0\) and
\(H(\lambda Y)\ne0\) for all \(\lambda\in\mathbb F_5\).
This condition tests the three pencils using the invariant coordinates.

Vertex relabeling and Frobenius also act on pairs \((B,L)\) of a
direction and a common line. Projection to \(B\) commutes with both
actions, giving bijections of the line fibers over conjugate directions.
On the free labeling locus, if \(F^e B=gB\) with \(g\in S_3\),
transport around that quotient orbit sends a line over \(B\) to
\(g^{-1}F^eL\). A choice returns to the same line exactly when this
map fixes its two-dimensional subspace. The line equations yield the
global exclusion in Theorem~\ref{thm:fixed-face-quadratic-empty}.
